\section{2017年高考题}
\subsection{2017年全国I卷（理）}\label{gk:2017年全国I卷（理）}


\clearpage


\subsection{2017年全国I卷（文）}\label{gk:2017年全国I卷（文）}


\clearpage


\subsection{2017年全国II卷（理）}\label{gk:2017年全国II卷（理）}


\clearpage


\subsection{2017年全国II卷（文）}\label{gk:2017年全国II卷（文）}


\clearpage


\subsection{2017年全国III卷（理）}\label{gk:2017年全国III卷（理）}


\clearpage


\subsection{2017年全国III卷（文）}\label{gk:2017年全国III卷（文）}


\clearpage


\subsection{2017年北京卷（理）}\label{gk:2017年北京卷（理）}

\begin{description}
    \item[一、]
    \begin{enumerate}
        \item 1
        \item 2
        \item 3
        \item 4
        \item 5
        \item 6
        \item 7
        \item 8
    \end{enumerate}
    \item[二、]
    \begin{enumerate}
      \addtocounter{enumi}{+8}   
        \item 1
        \item 2
        \item 3
        \item 4
        \item 1
        \item 6
    \end{enumerate}
    \item[三、]   
    \begin{enumerate}
      \addtocounter{enumi}{+14}   
        \item 1
        \item 2
        \item 3
        \item  已知抛物线 $C: y^{2} = 2 p x$ 过点 $P(1,1)$。 过点 $(0, \frac{1}{2})$ 作直线 $l$ 与抛物线 $C$ 交于不同的两点 $M, N$，过点 $M$ 作 $x$ 轴的垂线分别与直线 $OP$、$ON$ 交于点 $A, B$，其中 $O$ 为原点。

\begin{enumerate}[label=(\arabic*)]
    \item 求抛物线 $C$ 的方程，并求其焦点坐标和准线方程；
    \item 求证：$A$ 为线段 $BM$ 的中点。
\end{enumerate}
        \item  已知函数 $f(x) = e^{x} \cos x - x$。

\begin{enumerate}[label=(\arabic*)]
    \item 求曲线 $y = f(x)$ 在点 $(0, f(0))$ 处的切线方程；
    \item 求函数 $f(x)$ 在区间 $\left[0, \frac{\pi}{2}\right]$ 上的最大值和最小值。
\end{enumerate}
        \item  已知函数 $f(x) = e^{x} \cos x - x$。

\begin{enumerate}[label=(\arabic*)]
    \item 求曲线 $y = f(x)$ 在点 $(0, f(0))$ 处的切线方程；
    \item 求函数 $f(x)$ 在区间 $\left[0, \frac{\pi}{2}\right]$ 上的最大值和最小值。
\end{enumerate}
    \end{enumerate}
\end{description}
\clearpage


\subsection{2017年北京卷（文）}\label{gk:2017年北京卷（文）}

\begin{description}
    \item[一、]
    \begin{enumerate}
        \item 1
        \item 2
        \item 3
        \item 4
        \item 5
        \item 6
        \item 7
        \item 8
    \end{enumerate}
    \item[二、]
    \begin{enumerate}
      \addtocounter{enumi}{+8}   
        \item 1
        \item 2
        \item 3
        \item 4
        \item 1
        \item 6
    \end{enumerate}
    \item[三、]   
    \begin{enumerate}
      \addtocounter{enumi}{+14}   
        \item 1
        \item 2
        \item 3
        \item 
        \item 已知椭圆 $C$ 的两个顶点分别为 $A(-2,0)$, $B(2,0)$，焦点在 $x$ 轴上，离心率为 $\frac{\sqrt{3}}{2}$。

\begin{enumerate}[label=(\arabic*)]
    \item 求椭圆 $C$ 的方程；
    \item 点 $D$ 为 $x$ 轴上一点，过 $D$ 作 $x$ 轴的垂线交椭圆 $C$ 于不同的两点 $M, N$，过 $D$ 作 $AM$ 的垂线交 $BN$ 于点 $E$。求证：$\triangle BDE$ 与 $\triangle BDN$ 的面积之比为 $4:5$。
\end{enumerate}

        \item 已知函数 $f(x) = e^{x} \cos x - x$。

\begin{enumerate}[label=(\arabic*)]
    \item 求曲线 $y = f(x)$ 在点 $(0, f(0))$ 处的切线方程；
    \item 求函数 $f(x)$ 在区间 $\left[0, \frac{\pi}{2}\right]$ 上的最大值和最小值。
\end{enumerate}
    \end{enumerate}
\end{description}
\clearpage


\subsection{2017年江苏卷}\label{gk:2017年江苏卷}


\clearpage


\subsection{2017年上海卷（理）}\label{gk:2017年上海卷（理）}


\clearpage


\subsection{2017年上海卷（文）}\label{gk:2017年上海卷（文）}


\clearpage


\subsection{2017年山东卷（理）}\label{gk:2017年山东卷（理）}


\clearpage


\subsection{2017年山东卷（文）}\label{gk:2017年山东卷（文）}


\clearpage


\subsection{2017年天津卷（理）}\label{gk:2017年天津卷（理）}


\clearpage


\subsection{2017年天津卷（文）}\label{gk:2017年天津卷（文）}


\clearpage


\subsection{2017年浙江卷}\label{gk:2017年浙江卷}


\clearpage

